\clearpage
\VerdierPhysicalPageAnchor{0137}%
\begingroup
\small
\setlength{\emergencystretch}{1.5em}%
\setlength{\abovedisplayskip}{1pt plus 1pt minus 1pt}%
\setlength{\belowdisplayskip}{1pt plus 1pt minus 1pt}%

\VerdierSemanticAnchor{ch2-2-2-2-sm2-completion-diagram}
\noindent 补成如下交换图：

\begin{center}
  \begin{tikzpicture}[>=stealth,line width=.38pt,x=1.30cm,y=1.45cm,
    scale=.92,transform shape]
    \node (Zpp) at (0,2.20) {$Z''$};
    \node (Xp)  at (-1.25,1.10) {$X'$};
    \node (Yp)  at ( 1.25,1.10) {$Y'$};
    \node (Y)   at (0,0) {$Y$};
    \draw[->] (Zpp) -- node[above left=-1pt] {$t'\in S$} (Xp);
    \draw[->] (Zpp) -- node[above right=-1pt] {$m'$} (Yp);
    \draw[->] (Xp)  -- node[below left=-1pt] {$m$} (Y);
    \draw[->] (Yp)  -- node[below right=-1pt] {$t\in S$} (Y);
    \node at (2.50,-.05) {$.$};
  \end{tikzpicture}
\end{center}

\VerdierSemanticAnchor{ch2-2-2-2-composite-roof-d3}
\noindent 于是得到图式
$D_3=X\xleftarrow{st'}Z''\xrightarrow{pm'}Z$，其中复合态射 $st'$ 属于
$S$ (SM1)。
\VerdierSemanticAnchor{ch2-2-2-2-composite-class-independence}%
可以验证，$D_3$ 在
$\operatorname{Hom}_{\mathcal C(S^{-1})}(X,Z)$ 中的类
$\overline D_3$ 只依赖于类 $\overline D_1$ 与 $\overline D_2$。
\VerdierSemanticAnchor{ch2-2-2-2-composition-conclusion}%
态射 $\overline D_3$ 就是 $\mathcal C(S^{-1})$ 中态射
$\overline D_1$ 与 $\overline D_2$ 的复合。

\VerdierSemanticAnchor{ch2-2-2-3}
\smallskip
\noindent\textbf{注 2.2.3.}
\VerdierSemanticAnchor{ch2-2-2-3-a}%
a) 由定理 (2.2.2) 的第四个断言可知，配备函子
$Q\colon\mathcal C\longrightarrow\mathcal C(S^{-1})$ 的范畴
$\mathcal C(S^{-1})$ 是一个泛问题的解：
\VerdierSemanticAnchor{ch2-2-2-3-a-universal-problem}%
函子 $Q$ 把 $S$ 中的态射送到同构；而从 $\mathcal C$ 到某个范畴
$\mathcal A$、并把 $S$ 中态射送到同构的每个函子，都唯一地通过
$Q$ 分解。
\VerdierSemanticAnchor{ch2-2-2-3-a-unique-isomorphism}%
因此，由这一性质，配备函子 $Q$ 的范畴 $\mathcal C(S^{-1})$ 在唯一
同构意义下确定。

\VerdierSemanticAnchor{ch2-2-2-3-b}
\par\noindent b) 设 $\mathcal C$ 为范畴，$S$ 为 $\mathcal C$ 的一个态射
集合，不要求它是 $\mathcal C$ 的乘法系。
\VerdierSemanticAnchor{ch2-2-2-3-b-general-localization}%
总存在范畴 $\mathcal C(S^{-1})$ 以及把 $S$ 中态射送到同构的函子
$Q\colon\mathcal C\longrightarrow\mathcal C(S^{-1})$，使得从
$\mathcal C$ 到某个范畴 $\mathcal A$、并把 $S$ 中态射送到同构的
每个函子，都唯一地通过 $Q$ 分解。
\VerdierSemanticAnchor{ch2-2-2-3-b-multiplicative-comparison}%
当 $S$ 是 $\mathcal C$ 的乘法系时，定理 (2.2.2) 给出范畴
$\mathcal C(S^{-1})$ 的一个简单描述；它还给出函子 $Q$ 的若干附加
性质（断言 b) 与 c)），这些性质在一般情形下并不成立。

\VerdierSemanticAnchor{ch2-2-2-4}
\smallskip
\noindent\textbf{推论 2.2.4.}\quad
\begingroup\itshape
设 $\mathcal C$ 为范畴，$S$ 为 $\mathcal C$ 的乘法系。

\VerdierSemanticAnchor{ch2-2-2-4-a}
\par\noindent a) 集合 $S$ 是 $\mathcal C$ 的反范畴
$\mathcal C^\circ$ 的乘法系。
\VerdierSemanticAnchor{ch2-2-2-4-a-opposite-system}%
存在范畴的典范同构：

\VerdierSemanticAnchor{ch2-2-2-4-a-canonical-isomorphism}
\[
  \bigl(\mathcal C^\circ(S^{-1})\bigr)^\circ
  \longrightarrow \mathcal C(S^{-1})\,.
\]

\VerdierSemanticAnchor{ch2-2-2-4-a-hom-colimit-isomorphism}
\noindent 特别地，该同构诱导同构：
\[
  \operatorname{Hom}_{\mathcal C(S^{-1})}(X,Y)
  \longrightarrow
  \varinjlim_{Y\backslash S}\operatorname{Hom}_{\mathcal C}(X,\mathord\cdot)\,.
\]
\endgroup

\VerdierSemanticAnchor{ch2-2-2-4-p138-boundary}
\endgroup
